\documentclass[10pt,a4paper,fontset=none]{ctexart}
\usepackage[margin=1.65cm]{geometry}
\usepackage{xcolor}
\usepackage{listings}
\usepackage{longtable,booktabs,array}
\usepackage{enumitem}
\usepackage{hyperref}
\usepackage{fancyhdr}
\usepackage{titlesec}
\usepackage{tcolorbox}
\tcbuselibrary{breakable}
\usepackage{fontspec}
\setCJKmainfont{Songti SC}
\setCJKsansfont{Songti SC}
\setCJKmonofont{Songti SC}
\setmonofont{Menlo}[Scale=0.82]

\hypersetup{colorlinks=true, linkcolor=blue!45!black, urlcolor=blue!60!black}
\setlength{\headheight}{14pt}
\pagestyle{fancy}
\fancyhf{}
\lhead{DSA 期末考前学案}
\rhead{校订详细版}
\cfoot{\thepage}
\titleformat{\section}{\Large\bfseries\color{blue!45!black}}{\thesection}{0.7em}{}
\titleformat{\subsection}{\large\bfseries\color{black}}{\thesubsection}{0.6em}{}
\titleformat{\subsubsection}{\normalsize\bfseries\color{blue!35!black}}{\thesubsubsection}{0.5em}{}
\setlist[itemize]{itemsep=0.15em, topsep=0.22em, leftmargin=1.65em}
\setlist[enumerate]{itemsep=0.15em, topsep=0.22em, leftmargin=1.9em}
\newcolumntype{L}[1]{>{\raggedright\arraybackslash}p{#1}}

\definecolor{codebg}{RGB}{248,248,248}
\definecolor{darkgreen}{RGB}{0,100,0}
\lstdefinestyle{py}{
  language=Python,
  backgroundcolor=\color{codebg},
  basicstyle=\ttfamily\scriptsize,
  keywordstyle=\color{blue!70!black}\bfseries,
  commentstyle=\color{darkgreen},
  stringstyle=\color{red!55!black},
  showstringspaces=false,
  frame=single,
  rulecolor=\color{black!20},
  breaklines=true,
  tabsize=4,
  columns=fullflexible,
  keepspaces=true
}
\lstdefinestyle{plain}{
  backgroundcolor=\color{codebg},
  basicstyle=\ttfamily\scriptsize,
  frame=single,
  rulecolor=\color{black!15},
  breaklines=true,
  columns=fullflexible,
  keepspaces=true
}

\newtcolorbox{sourcebox}{colback=blue!3!white,colframe=blue!40!black,title=资料依据,breakable}
\newtcolorbox{methodbox}{colback=green!2!white,colframe=green!40!black,title=本题实际代码方法,breakable}
\newtcolorbox{warnbox}{colback=red!2!white,colframe=red!45!black,title=需要警惕,breakable}
\newtcolorbox{checklistbox}{colback=yellow!5!white,colframe=yellow!45!black,title=落地检查,breakable}

\title{\bfseries DSA 期末考前学案：校订详细版\\\large 严格依据错题 PDF、对应 Python 代码、tree.md 与 graph.md}
\author{根据当前文件夹资料重新校订}
\date{2026 年 6 月}

\begin{document}
\maketitle
\tableofcontents
\newpage

\section*{校订说明}
\addcontentsline{toc}{section}{校订说明}
这份版本专门修正上一版的问题：不再把没有出现在 \texttt{tree.md} 或 \texttt{graph.md} 中的内容说成“md 里已有”，也不再把 PDF 中你实际采用的代码方法改写成另一套更短或更常见的做法。每道题都按三件事来写：
\begin{enumerate}
  \item \textbf{md 中确实相关的内容：}只写两个 md 文件里实际出现的模板、概念或例子。
  \item \textbf{PDF/对应 Python 代码的实际做法：}以当前 PDF 的 Listing 和同名 \texttt{.py} 文件为准。
  \item \textbf{从模板到本题的迁移：}解释你的代码为什么这样改，而不是另起炉灶。
\end{enumerate}

\begin{warnbox}
特别校订的几处：\texttt{实现堆结构.py} 实际是手写 \texttt{BinHeap}，不是直接用 \texttt{heapq}；\texttt{兔子与樱花.py} 实际是图对象 + 无优先队列 Dijkstra，不是 Floyd；\texttt{丛林中的路.py} 实际是邻接矩阵 + Prim，不是 Kruskal；\texttt{括号嵌套二叉树.py} 实际是按下标递归解析，不是“最外层逗号 split”。这些都会按实际代码讲。
\end{warnbox}

\section{两个 md 文件实际覆盖了什么}

\subsection{tree.md 的真实内容}
\begin{sourcebox}
\texttt{tree.md} 实际覆盖：
\begin{itemize}
  \item 二叉树节点类、前序/中序/后序递归遍历。
  \item 二叉树广度优先遍历。
  \item 从层次遍历建树。
  \item 普通树按“左儿子右兄弟”转换成二叉树。
  \item 前序 + 中序重建二叉树。
  \item 非递归前序/中序遍历。
  \item 二叉堆 ADT、手写堆的上浮/下沉/buildHeap、堆排序。
  \item Python \texttt{heapq} 的最小堆用法。
  \item BST 的节点、插入、查找、删除。
  \item Huffman/最优二叉树：最小 WPL、真正构造 Huffman 树、Huffman 解码、木板切割最小费用。
\end{itemize}
\end{sourcebox}

它没有直接覆盖的作业细节包括：括号嵌套二叉树的按下标解析、扩展二叉树的 \texttt{.} 空节点输入、BST 合法性上下界验证、完全二叉树编号区间计数。这些属于作业/Python 文件中的具体题型拓展。

\subsection{graph.md 的真实内容}
\begin{sourcebox}
\texttt{graph.md} 实际覆盖：
\begin{itemize}
  \item 图的基本概念、邻接矩阵、邻接列表、Vertex/Graph 类。
  \item 词梯问题：桶建图 + BFS。
  \item BFS 原理、距离/前驱/颜色、路径回溯。
  \item 数轴 BFS“抓住那头牛”。
  \item 普通迷宫 BFS，带 parent 的路径恢复。
  \item 迷宫变形：剩余消除障碍能力作为状态维度。
  \item 进入守卫等不同代价迷宫：普通 BFS 失效，改优先队列/Dijkstra。
  \item 更复杂钥匙/守卫状态设计的讨论。
  \item DFS、骑士周游、城堡问题、踩方格、算 24。
  \item 拓扑排序、AOE 关键路径、强连通分支。
  \item Dijkstra，包括优先队列版本和不用优先队列版本。
  \item Floyd，包括 \texttt{dist} 和 \texttt{prev}。
  \item Prim，包括对象图版本、\texttt{heapq} 版本、邻接矩阵版本。
  \item Kruskal 的文字定义。
\end{itemize}
\end{sourcebox}

它没有直接覆盖的作业细节包括：A Walk Through the Forest 的“从 2 跑 Dijkstra 后按距离下降 DFS 计数”、兔子与樱花中你实际用“每次查询重新跑无优先队列 Dijkstra”、连接所有点的最小费用中你实际先建完全图再 Prim。这些必须以 Python 文件为准讲解。

\section{总体薄弱点：按实际代码归纳}

\subsection{薄弱点一：模板和题目输入之间的转换}
你的很多代码是从 md 模板改出来的，但题目输入形式发生了变化：
\begin{itemize}
  \item \texttt{tree.md} 中前序 + 中序重建是模板；作业 22 需要多组 EOF，并在建树后后序输出。
  \item \texttt{tree.md} 中层次建树是模板；作业 54 用它先建树，再做 BST 上下界验证。
  \item \texttt{graph.md} 中词梯是从文件读单词；作业“单词序列”从一行输入读 \texttt{wordlist}，还要把 start/end 加进去。
  \item \texttt{graph.md} 中 Dijkstra 是对象图模板；兔子与樱花直接使用这个对象图结构，但每次查询前要 reset。
\end{itemize}

\subsection{薄弱点二：相似题名下算法目标不同}
不要只看题目像不像，而要看答案目标：
\begin{itemize}
  \item 仙岛求药：普通 BFS，目标是最短步数。
  \item 打怪救公主：不是最短步数，而是最大剩余血量，所以用 \texttt{max\_blood} 二维数组。
  \item 拯救行动：最短时间，但边权 1/2 不同，所以用 Dijkstra。
  \item A Walk Through the Forest：最终是路线数量，不是最短路长度；Dijkstra 只是中间步骤。
\end{itemize}

\subsection{薄弱点三：手写结构和库函数的对应关系}
本次代码里既有手写结构，也有库函数：
\begin{itemize}
  \item 实现堆结构：手写 \texttt{BinHeap}，与 \texttt{tree.md} 的手写堆最接近。
  \item Huffman：使用 \texttt{heapq}，与 \texttt{tree.md} 的最小 WPL 代码最接近。
  \item 拓扑排序、拯救行动、Prim：都使用 \texttt{heapq}，但含义不同。拓扑堆存点编号，Dijkstra 堆存距离，Prim 堆存连接边权。
\end{itemize}

\section{树类题逐题校订复盘}

\subsection{22 根据二叉树前中序序列建树}
\begin{sourcebox}
md 依据：\texttt{tree.md} 中“前序 + 中序重建二叉树”的基本原则和 \texttt{Buildtree(preorder, inorder)} 代码。
\end{sourcebox}
\begin{methodbox}
你的实际代码：定义 \texttt{Treenode}；用 \texttt{Buildtree(preorder, inorder)} 递归建树；再用 \texttt{postorder(root, ans)} 输出后序；主程序用 \texttt{while True + try/except EOFError} 处理多组输入。

代码建立思路：这题最直接依赖 \texttt{tree.md} 里的“前序 + 中序重建二叉树”模板。现成部分是：前序第一个元素当根，在中序中找到根的位置，并据此切出左、右子树。你新写/修改的部分是两个：第一，建好树后接上后序遍历函数，因为题目要输出后序；第二，主程序改成 EOF 多组输入。组合方式就是：读一组前序/中序 $\rightarrow$ 用 md 模板建树 $\rightarrow$ 后序遍历收集字符 $\rightarrow$ 输出。
\end{methodbox}
迁移关系：
\begin{enumerate}
  \item md 里已经有前序根 \texttt{preorder[0]} 和中序切左右。
  \item 本题多了“输出后序”，所以不是建完就结束，而是后序遍历收集答案。
  \item 本题多组到 EOF，所以主循环不是只读一次。
\end{enumerate}
易错点：前序左子树切片必须用左子树长度 \texttt{l=len(ino\_left)}，你的代码正是 \texttt{preorder[1:1+l]}。

\subsection{23 括号嵌套二叉树}
\begin{sourcebox}
md 依据：只有二叉树节点和前序/中序遍历模板；\texttt{tree.md} 没有直接给括号嵌套解析模板。
\end{sourcebox}
\begin{methodbox}
你的实际代码：\texttt{buildtree(s, i)} 按字符串下标递归解析。遇到 \texttt{*} 返回空和新下标；读到根字符后，如果后面是 \texttt{(}，递归读左子树，遇逗号后递归读右子树，遇 \texttt{)} 返回。

代码建立思路：这题只借 \texttt{tree.md} 的二叉树节点类和递归遍历思想，解析括号串这一块是本题新写的。你先保留“一个节点有 \texttt{left/right}”的树结构，再把输入字符串看成递归语法：当前字符是根，\texttt{*} 是空树，括号里依次是左子树和右子树。新写的关键是让 \texttt{buildtree} 返回“子树根 + 读到的新下标”，这样递归解析和字符串位置同步推进；建完后再接上 \texttt{tree.md} 里熟悉的前序/中序输出。
\end{methodbox}
这题最重要的是“函数返回两个东西”：子树根节点和当前解析到的新下标。这样避免了自己手动找最外层逗号。
\begin{lstlisting}[style=py]
root.left, i = buildtree(s, i)
if i < len(s) and s[i] == ",":
    i += 1
    root.right, i = buildtree(s, i)
\end{lstlisting}
复习时要把它归到“树遍历 + 字符串递归解析”，不要误以为 md 已经有完整解析模板。

\subsection{24 二叉树}
\begin{sourcebox}
md 依据：\texttt{tree.md} 中完全二叉树编号性质：左孩子 \texttt{2p}，右孩子 \texttt{2p+1}。
\end{sourcebox}
\begin{methodbox}
你的实际代码：不建树。读入 \texttt{n,m}，其中 \texttt{n} 是子树根编号，\texttt{m} 是总编号上界；用 \texttt{left/right} 表示当前层在这棵子树中的编号区间，逐层扩张并截断。

代码建立思路：这题依赖 \texttt{tree.md} 中完全二叉树的编号关系，而不是依赖节点建树模板。现成概念是：编号为 $p$ 的节点，左孩子是 $2p$，右孩子是 $2p+1$。你新写的做法是把“一个子树在每一层上的所有编号”压缩成区间 \texttt{[left,right]}：下一层区间变为 \texttt{[left*2,right*2+1]}。每层只统计不超过 \texttt{m} 的编号个数，所以整题组合为：根编号区间初始化 $\rightarrow$ 按完全树编号扩层 $\rightarrow$ 每层截断计数。
\end{methodbox}
关键公式：
\begin{lstlisting}[style=py]
ans += min(m, right) - left + 1
right = right * 2 + 1
left = left * 2
\end{lstlisting}
这是从 md 的“编号性质”进一步变成“区间计数”。这里不是 md 里直接给出的题型，但完全依赖那个性质。

\subsection{25 二叉搜索树的层次遍历}
\begin{sourcebox}
md 依据：BST 插入、二叉树 BFS/层次遍历。
\end{sourcebox}
\begin{methodbox}
你的实际代码：手写 \texttt{TreeNode} 和 \texttt{BinarySearchTree.put}；插入时小于走左，大于走右，相等不处理；最后 \texttt{bfs(root, ans)} 用列表队列做层次遍历。

代码建立思路：这题是 \texttt{tree.md} 里 BST 插入模板和二叉树层次遍历模板的组合。第一段沿用 BST 的比较逻辑：从根开始，小于当前节点就往左，大于当前节点就往右，直到空位置插入；作业中特别修改的是相等元素不插入。第二段沿用层次遍历：把根放入队列，弹出节点就记录值，并把非空左右孩子入队。组合方式就是：输入序列逐个插入成 BST $\rightarrow$ 对生成的 BST 做 BFS $\rightarrow$ 输出层序结果。
\end{methodbox}
本题从 md 到作业的变化是：不是单独考 BST，也不是单独考 BFS，而是“先用输入序列构造 BST，再按层次输出”。重复值忽略是作业题面特有点，md 的通用 BST 插入代码没有强调这个输出要求。

\subsection{26 实现堆结构}
\begin{sourcebox}
md 依据：\texttt{tree.md} 中手写 \texttt{BinHeap} 的初始化、\texttt{percUp}、\texttt{percDown}、\texttt{delMin}。
\end{sourcebox}
\begin{methodbox}
你的实际代码：使用手写 \texttt{BinHeap}，不是 \texttt{heapq}。每组测试新建一个堆，操作 1 调 \texttt{insert}，操作 2 调 \texttt{delMin} 并打印。

代码建立思路：这题直接依赖 \texttt{tree.md} 的手写二叉堆 ADT，而不是 Python 库堆。现成部分是用数组保存完全二叉树、用下标关系找父子节点、插入后 \texttt{percUp}、删除堆顶后 \texttt{percDown}。你新写/修改的部分是把比较方向统一成小顶堆，并把堆操作包进题目的操作流：操作 1 读一个数插入，操作 2 删除并输出当前最小值。组合方式就是：每组数据初始化空堆 $\rightarrow$ 按操作调用堆方法 $\rightarrow$ 删除时打印返回值。
\end{methodbox}
需要特别注意：\texttt{tree.md} 里手写 \texttt{percUp} 的比较方向写成了大顶堆倾向，而你的作业代码改成 \texttt{<}，实现的是最小堆。这是从笔记到题目要求的关键修正。
\begin{lstlisting}[style=py]
if self.heaplist[i] < self.heaplist[i//2]:
    self.heaplist[i], self.heaplist[i//2] = ...
\end{lstlisting}

\subsection{27 树的转换}
\begin{sourcebox}
md 依据：普通树转左儿子右兄弟二叉树的 \texttt{convert(root)} 模板。
\end{sourcebox}
\begin{methodbox}
你的实际代码：先根据 \texttt{d/u} 字符串真正建一棵普通树，节点有 \texttt{children} 和 \texttt{parent}；再调用 \texttt{convert} 生成二叉树；最后分别用 \texttt{length(root)} 和 \texttt{height(b\_root)} 输出原树高度和二叉树高度。

代码建立思路：这题依赖 \texttt{tree.md} 的“普通树转左儿子右兄弟二叉树”模板，但题目没有直接给一棵现成的普通树，所以你先新写了建普通树的过程。读到 \texttt{d} 就创建新孩子并下移，读到 \texttt{u} 就回到父节点；这样得到带 \texttt{children} 列表的普通树。接着套 md 的转换思想：第一个孩子变成二叉树左孩子，后续兄弟沿右指针串起来。最后再分别写两个高度函数，原树高度按 \texttt{children} 递归，转换后的二叉树高度按 \texttt{left/right} 递归。
\end{methodbox}
这题不是“只模拟高度”的写法，而是与 md 中转换模板高度一致：先建普通树，再转换，再递归求高度。注意你的 \texttt{height(None)} 返回 \texttt{-1}，所以叶子二叉树高度为 0；\texttt{length} 也按边数口径处理。

\subsection{28 Huffman 编码树}
\begin{sourcebox}
md 依据：\texttt{tree.md} 中“计算最小 WPL”的 \texttt{heapq} 版本和最优二叉树构造思想。
\end{sourcebox}
\begin{methodbox}
你的实际代码：\texttt{optimal\_binary\_tree(weights)} 直接 \texttt{heapq.heapify}，每次弹出两个最小值，合并值加入 \texttt{total\_wpl}，再压回堆。主程序按测试组读取。

代码建立思路：这题依赖 \texttt{tree.md} 中 Huffman/最优二叉树的“最小 WPL”代码，而不是依赖真正输出编码的树结构。现成部分是 \texttt{heapq} 小根堆：每次取出两个最小权值合并，合并代价加入总 WPL，再把新权值压回。你新写/修改的部分主要是读多组测试数据，以及把函数命名为 \texttt{optimal\_binary\_tree} 只返回最小带权路径长度。组合方式就是：权值列表建堆 $\rightarrow$ 反复两两合并 $\rightarrow$ 累加所有合并代价。
\end{methodbox}
这里和 md 的最小 WPL 代码基本一致。不要把它讲成“必须真正构造 Huffman 树”；你的作业代码只计算 WPL。

\subsection{49 二叉树的深度}
\begin{sourcebox}
md 依据：二叉树节点和递归高度思想。
\end{sourcebox}
\begin{methodbox}
你的实际代码：先创建 \texttt{nodes=[None]*(n+1)}，为 1 到 n 每个编号建 \texttt{BinaryNode}；再读每个节点的左右孩子编号并连接对象；最后 \texttt{height(nodes[1])} 输出深度。

代码建立思路：这题依赖 \texttt{tree.md} 的二叉树节点类和递归求高度思想。现成部分是“节点有左孩子、右孩子，树高由左右子树高度递归得到”；你新写的部分是把题目的编号输入转换成真实节点连接。代码先为每个编号预建对象，再根据输入的左右孩子编号把 \texttt{left/right} 指针接好，编号 \texttt{-1} 表示空。组合方式就是：编号表建节点 $\rightarrow$ 按输入连边 $\rightarrow$ 从 1 号根递归算高度/深度。
\end{methodbox}
和更短的数组递归不同，你的代码走的是“编号输入转节点对象”的路线。这与 \texttt{tree.md} 中用类表示树节点的思路一致。

\subsection{50 扩展二叉树}
\begin{sourcebox}
md 依据：二叉树节点、递归遍历。md 没有直接给扩展先序 \texttt{.} 输入模板。
\end{sourcebox}
\begin{methodbox}
你的实际代码：\texttt{build\_tree(seq)} 内部定义 \texttt{helper()}，用 \texttt{nonlocal idx} 顺序读先序串；遇到 \texttt{.} 返回 \texttt{None}；否则建节点并递归建左、右子树。最后输出中序和后序。

代码建立思路：这题依赖 \texttt{tree.md} 的二叉树递归遍历，但建树方式是作业新加的“扩展先序”。现成部分是节点结构和中序/后序遍历函数；新写部分是 \texttt{helper()} 顺序消费字符串：每读到一个普通字符就建节点，然后递归读它的左子树和右子树；每读到 \texttt{.} 就返回空节点。组合方式就是：扩展先序串 $\rightarrow$ 用全局/闭包下标递归建树 $\rightarrow$ 对建出的树分别中序、后序输出。
\end{methodbox}
本题核心是“带空标记的先序能唯一建树”。从 md 的遍历模板迁移时，新增的是顺序读入下标和空节点出口。

\subsection{53 根据二叉树中后序序列建树}
\begin{sourcebox}
md 依据：前序 + 中序重建的切分思想。md 没有直接给中序 + 后序代码。
\end{sourcebox}
\begin{methodbox}
你的实际代码：\texttt{Buildtree(postorder, inorder)}；根是 \texttt{postorder[-1]}；用中序切左右；按左子树长度切后序；建树后前序输出。

代码建立思路：这题依赖 \texttt{tree.md} 中“用中序确定左右子树范围”的重建思想，但把前序模板改成了后序版本。现成逻辑是：必须在中序里找到根，根左边是左子树，根右边是右子树。你修改的关键是根不再取 \texttt{preorder[0]}，而是取 \texttt{postorder[-1]}；切后序时要排除最后的根，并按左子树长度拆成左右两段。组合方式就是：中序定位根 $\rightarrow$ 切左右子树序列 $\rightarrow$ 递归建树 $\rightarrow$ 前序遍历输出。
\end{methodbox}
这是第 22 题的对称迁移：根的位置从前序开头换成后序末尾，输出方式从后序换成前序。

\subsection{54 验证二叉搜索树}
\begin{sourcebox}
md 依据：层次遍历建树模板、BST 的基本性质。md 没有直接给“验证 BST 上下界递归”。
\end{sourcebox}
\begin{methodbox}
你的实际代码：先用 \texttt{build\_tree\_level(level\_list)} 从层序输入建树，\texttt{\#} 表示空；再用 \texttt{is\_bst(node, low, high)} 递归验证，每个节点必须满足 \texttt{low < val < high}。

代码建立思路：这题分两段，第一段依赖 \texttt{tree.md} 的层次遍历建树模板：用队列按顺序给当前节点接左、右孩子，遇到 \texttt{\#} 就不建节点。第二段依赖 BST 的性质，但验证函数是本题新写的增强版：不能只比较父子大小，而是把祖先给出的上下界 \texttt{low/high} 一路传下去。组合方式就是：层序字符串建出普通二叉树 $\rightarrow$ 从根开始做上下界递归 $\rightarrow$ 左子树收紧上界，右子树收紧下界。
\end{methodbox}
这题的拓展点在于：不能只看一个节点和左右孩子的大小关系，而要把祖先限制一路传下去。你的代码已经用上下界处理了这个问题。

\section{图类题逐题校订复盘}

\subsection{1 单词序列}
\begin{sourcebox}
md 依据：\texttt{graph.md} 词梯问题中的桶建图、Vertex/Graph 类、BFS 的 color/distance/pred。
\end{sourcebox}
\begin{methodbox}
你的实际代码：几乎沿用 md 的对象图结构；\texttt{buildGraph(wordlist)} 用通配桶把只差一个字母的单词连边；BFS 用 \texttt{deque}，设置 color、distance、pred；答案输出 \texttt{e.getDistance()+1}，不可达输出 0。

代码建立思路：先直接拿 \texttt{graph.md} 里的 \texttt{Vertex/Graph} 邻接表类，作为存单词图的骨架；再拿词梯的“通配桶建图”代码，把输入方式从读文件改成读一行 \texttt{wordlist}，并把起点、终点补进词表。建完图后接上 md 的 BFS 模板，用 \texttt{distance} 表示从起点到某个单词需要换几次；最后新写的部分主要是输入输出和答案转换：如果终点没被访问输出 0，否则输出边数距离加 1，表示序列中单词个数。
\end{methodbox}
从 md 到作业的具体改动：
\begin{itemize}
  \item md 是从文件读单词，本题从一行 \texttt{wordlist=input().split()} 读。
  \item 你的代码把 start 和 end 也 append 到 wordlist，避免起终点不在字典中。
  \item 输出的是单词序列长度，所以边数距离要加 1。
\end{itemize}

\subsection{2 仙岛求药}
\begin{sourcebox}
md 依据：普通迷宫 BFS，二维 \texttt{visited}，四方向扩展。
\end{sourcebox}
\begin{methodbox}
你的实际代码：\texttt{bfs(maze,M,N,sx,sy)} 从 \texttt{@} 出发，队列元素为 \texttt{(x,y,dis)}；遇到 \texttt{*} 返回距离；\texttt{\#} 不可走；\texttt{0 0} 结束。

代码建立思路：这里直接用 \texttt{graph.md} 的普通迷宫 BFS 模板，不需要建显式 Graph 类，因为网格本身就是隐式图。保留模板中的队列、四方向数组、二维 \texttt{visited}；新写/修改的是读多组输入、扫描 \texttt{@} 起点、把目标判断改成 \texttt{maze[nx][ny]=='*'}，并把墙判断固定为 \texttt{\#} 不能走。组合顺序就是：读迷宫 $\rightarrow$ 找起点 $\rightarrow$ BFS $\rightarrow$ 返回最先到达终点的步数。
\end{methodbox}
这题和 md 普通迷宫模板最接近。注意你的起始距离是 0，因此输出的是移动步数，不是包含起点的格子数。

\subsection{3 变换的迷宫}
\begin{sourcebox}
md 依据：迷宫变形问题中“同一个位置在不同状态下不同，要用三维 visited”的思想。
\end{sourcebox}
\begin{methodbox}
你的实际代码：\texttt{visited[M][N][K]}，状态是位置加时间模数；如果下一格是 \texttt{\#} 且 \texttt{(dis+1)\%K != 0}，则不能进入；到达 \texttt{E} 返回时间，不可达输出 \texttt{Oop!}。

代码建立思路：这题借的是 \texttt{graph.md} 里“状态 BFS”的搭法，而不是照搬剩余能力的含义。现成部分是队列 BFS、四方向移动、\texttt{visited} 升成三维；修改部分是把第三维从“剩余能力”换成“时间模数 \texttt{t\%K}”。新写的关键判断是：普通格可以按 BFS 进入，墙 \texttt{\#} 只有在下一时刻 \texttt{(dis+1)\%K==0} 才能通过。代码组合为：网格位置 + 当前时间共同定义状态，入队时同步更新 \texttt{dis+1} 和时间模数。
\end{methodbox}
这里与 md 的“剩余消除能力 k”不完全相同：md 的第三维是剩余能力，本题第三维是时间模数。但迁移原则一致：二维位置不够表示状态，必须升维。

\subsection{1 拓扑排序}
\begin{sourcebox}
md 依据：拓扑排序的入度数组 + 队列版本。
\end{sourcebox}
\begin{methodbox}
你的实际代码：用 \texttt{heapq} 维护当前入度为 0 的点，实现“编号小优先”；图数组 \texttt{G=[[] for \_ in range(n+1)]}，顶点从 1 到 n；输出格式为 \texttt{v1 v2 ...}。

代码建立思路：主体来自 \texttt{graph.md} 的拓扑排序：建邻接表、统计入度、不断取入度为 0 的点、删除它的出边。你修改的核心只有一个：把 md 里的普通队列换成 \texttt{heapq}，这样多个点同时入度为 0 时会先弹出编号最小者。其他新写部分主要是作业格式适配：顶点编号从 1 开始，读入边后更新 \texttt{indegree[v]}，最后把结果加上 \texttt{v} 前缀输出。
\end{methodbox}
从 md 到作业的变化：md 讲的是普通队列输出任一拓扑序；你的代码为了编号小优先，把队列换成了最小堆。

\subsection{2 丛林中的路}
\begin{sourcebox}
md 依据：Prim 的邻接矩阵版本，以及 MST 的定义。md 也有 Kruskal 文字，但你的代码没有用 Kruskal。
\end{sourcebox}
\begin{methodbox}
你的实际代码：\texttt{Graph\_adj\_mat} 存邻接矩阵；输入字母用 \texttt{ord(ch)-ord('A')} 转编号；无向边双向写入矩阵；调用 \texttt{prim(G,0)}，堆中存 \texttt{(边权, 顶点编号)}。

代码建立思路：这里使用的是 \texttt{graph.md} 的 Prim 邻接矩阵版本。现成部分是邻接矩阵图结构和 Prim 的“从已选集合向外找最小边”逻辑；新写/修改部分是输入解析，因为题目用大写字母表示点，还按每个点后面跟若干条边的形式给数据。你的代码先把字母转成 0 开始的编号，再把每条无向边同时写入 \texttt{G[u][v]} 和 \texttt{G[v][u]}，最后把矩阵图交给 \texttt{prim(G,0)} 汇总 MST 边权。
\end{methodbox}
这题必须校订：它不是 Kruskal。复习时应和 \texttt{graph.md} 中“基于邻接矩阵的 Prim”对应起来。

\subsection{3 A Walk Through the Forest}
\begin{sourcebox}
md 依据：Dijkstra 对象图模板、\texttt{heapq} 过期状态处理、DFS/记忆化思想可从图搜索部分迁移。md 没有直接给这道题的组合。
\end{sourcebox}
\begin{methodbox}
你的实际代码：建 1..n 的对象图，读无向带权边；从顶点 2 跑 Dijkstra，得到每个点到家的最短距离；然后 \texttt{dfs\_count\_routes(u,memo)} 从 1 开始计数，只能走向 \texttt{v.distance < u.distance} 的邻居。

代码建立思路：这题是组合题。第一段借 \texttt{graph.md} 的对象图 + Dijkstra 优先队列版：用 \texttt{Vertex/Graph} 存无向带权边，并从终点 2 跑最短路，得到“每个点离家有多近”。第二段不是 md 的原样代码，而是你新写的 DFS 计数：从 1 出发，只允许走向 \texttt{distance} 更小的邻居，因为题目要求每一步都更接近家。最后用 \texttt{memo} 记忆化每个点到 2 的合法路线数，把 Dijkstra 的距离结果和 DFS 的路线计数组合起来。
\end{methodbox}
这题要特别分清两层：
\begin{enumerate}
  \item Dijkstra 的答案不是最终输出，而是为“更接近家”提供判断依据。
  \item DFS 的终点是编号 2，\texttt{memo[u]} 记录从 u 到家的合法路线数。
\end{enumerate}

\subsection{4 打怪救公主}
\begin{sourcebox}
md 依据：普通迷宫 BFS、以及“状态中资源会影响后续”的思想。md 没有直接给最大剩余血量版本。
\end{sourcebox}
\begin{methodbox}
你的实际代码：\texttt{max\_blood[x][y]} 存到达每个格子的最大剩余血量；队列元素是 \texttt{(x,y,blood)}；进入数字格扣血，扣到 0 或以下不能继续；如果 \texttt{new\_blood > max\_blood[nx][ny]} 才入队。

代码建立思路：基础移动框架来自 \texttt{graph.md} 的普通迷宫 BFS：队列、四方向、边界判断、墙/可走格判断。你新写的关键是把 \texttt{visited} 从布尔数组改成 \texttt{max\_blood}，因为同一个格子“到过”不够，还要看剩余血量是否更优。每次扩展时先按格子内容计算 \texttt{new\_blood}，如果血量没有超过该格历史最好值就丢弃；如果更高，才更新数组并入队。也就是说：BFS 的壳保留，判重条件从“没访问过”改成“资源状态更好”。
\end{methodbox}
这题不是求最短路，所以不能用普通 \texttt{visited=True/False}。你的代码正确保留了“更高血量重访同一格”的可能。

\subsection{5 兔子与樱花}
\begin{sourcebox}
md 依据：\texttt{graph.md} 中 Dijkstra 不用优先队列版本、Vertex/Graph 类、前驱 \texttt{pred} 恢复路径。md 的 Floyd 部分不是你这题实际采用的方法。
\end{sourcebox}
\begin{methodbox}
你的实际代码：地点名映射到 Vertex；道路是无向边；每次查询时对起点运行 \texttt{dijkstra\_no\_pq}，通过 \texttt{pred} 从终点回溯路径，再按 \texttt{A->(w)->B} 输出。

代码建立思路：这题直接沿用 \texttt{graph.md} 的 \texttt{Vertex/Graph} 对象图，以及无优先队列 Dijkstra 的主体：每轮扫描未访问点中当前距离最小的点，再松弛它的邻居。你新写/修改的部分有三块：第一，把城市字符串映射成顶点对象，而不是用数字编号；第二，因为有多次查询，每次跑 Dijkstra 前都要 reset 所有顶点的 \texttt{distance/pred/visited}；第三，Dijkstra 结束后不是只输出距离，而是沿 \texttt{pred} 从终点回溯路径，再根据相邻城市查边权，格式化成题目要求的路径字符串。
\end{methodbox}
这题必须校订：它不是 Floyd。虽然 Floyd 也适合多查询，但你的代码实际选择了“每次查询跑一次无优先队列 Dijkstra”。讲解时应重点解释 reset、visited 集合、pred 路径恢复。

\subsection{56 拯救行动}
\begin{sourcebox}
md 依据：\texttt{graph.md} 中“进入守卫格代价不同，所以普通 BFS 转 Dijkstra”的迷宫例子。
\end{sourcebox}
\begin{methodbox}
你的实际代码：每组读迷宫，找到 \texttt{r} 和 \texttt{a}；\texttt{dist[x][y]} 存最短时间；堆中存 \texttt{(steps,x,y)}；进入 \texttt{x} 代价 2，否则代价 1；遇到终点打印 steps，不可达打印 \texttt{Impossible}。

代码建立思路：这题基本套 \texttt{graph.md} 的“不同代价迷宫 Dijkstra”。现成部分是 \texttt{heapq} 小根堆、\texttt{dist} 二维数组、每次弹出当前最短时间状态、四方向松弛。新写/修改的是题目字符含义：起点是 \texttt{r}，终点是 \texttt{a}，墙不能进，进入 \texttt{x} 的代价是 2，进入其他可走格代价是 1。组合方式就是把网格当作隐式带权图，不显式建边，用 Dijkstra 在扩展邻格时即时计算边权。
\end{methodbox}
这题和 md 的“解题思路二”高度一致。它最核心的判断是：普通 BFS 按步数扩展，但本题按总时间最小扩展。

\subsection{57 连接所有点的最小费用}
\begin{sourcebox}
md 依据：Prim 的 \texttt{heapq} 对象图版本、完全图概念、MST 定义。
\end{sourcebox}
\begin{methodbox}
你的实际代码：用 \texttt{ast.literal\_eval} 读 points；先为每对点计算曼哈顿距离并加入对象图，包含双向边；再调用 \texttt{prim(g,g.getVertex(0))}。Prim 中用 \texttt{inTree} 防止重复加入，用 \texttt{distance} 记录当前连接代价。

代码建立思路：这里用了 \texttt{graph.md} 的对象图 Prim/heapq Prim 思路。现成部分是 \texttt{Vertex/Graph} 存边、Prim 用堆维护“把某个新点接入当前树的最小代价”、用 \texttt{inTree} 跳过已经加入 MST 的点。你新写的部分是把点坐标转成完全图：两层循环枚举每一对点，计算曼哈顿距离，并作为无向边加入图。最后组合为：输入坐标 $\rightarrow$ 显式建立完全带权图 $\rightarrow$ 从 0 号点运行 Prim $\rightarrow$ 累加每次接入新点的最小边权。
\end{methodbox}
这里也要校订：你的代码不是“不建完全图边”的简化 Prim，而是先显式建了完全图对象图，再跑 Prim。学案应讲这个实际方法，同时提醒它与 md 的 Prim 对象图版本最相近。

\section{横向对比：md 模板如何被你实际改动}

\begin{longtable}{L{3.2cm}L{4.2cm}L{6.7cm}}
\toprule
\textbf{md 模板/概念} & \textbf{对应作业题} & \textbf{你实际做出的改动} \\
\midrule
前序+中序建树 & 22 & 保留 \texttt{Buildtree}，新增 EOF 多组输入和后序输出。\\
二叉树节点与遍历 & 23、50、49 & 23 增加按下标解析括号；50 增加 \texttt{.} 空节点和 \texttt{nonlocal idx}；49 从编号输入连接对象。\\
层次建树 & 54 & 用 \texttt{\#} 建空节点，再上下界验证 BST。\\
手写堆 & 26 & 比较方向改成小顶堆；按操作流插入/删除。\\
\texttt{heapq} Huffman & 28 & 基本照最小 WPL 模板，多组输入。\\
左儿子右兄弟转换 & 27 & 先按 \texttt{d/u} 建普通树，再调用 \texttt{convert}，最后求两种高度。\\
词梯桶建图+BFS & 单词序列 & 从输入列表建桶；把 start/end 加入词表；输出边数+1。\\
普通迷宫 BFS & 仙岛求药 & 起点 \texttt{@}，终点 \texttt{*}，墙 \texttt{\#}，0 0 结束。\\
状态 BFS & 变换的迷宫 & 第三维从“剩余能力”变成“时间模数”。\\
不同代价迷宫 Dijkstra & 拯救行动 & 完全对应，进入 \texttt{x} 代价 2。\\
最优资源状态思想 & 打怪救公主 & 用 \texttt{max\_blood} 而不是布尔 visited。\\
拓扑排序队列 & 拓扑排序 & 队列改最小堆，保证编号小优先。\\
Prim 邻接矩阵 & 丛林中的路 & 字母转编号，矩阵存边，堆式 Prim。\\
Dijkstra 对象图 & 兔子与樱花 & 每个查询 reset 后跑无优先队列 Dijkstra，再 pred 回溯。\\
Dijkstra + DFS 自行组合 & A Walk Through the Forest & 从 2 跑 Dijkstra，再按距离下降 DFS 计数。\\
Prim 对象图 & 连接所有点最小费用 & 先显式建完全图，再用对象图 Prim。\\
\bottomrule
\end{longtable}

\section{graph.md 中有具体代码的算法清单与题目映射}

\subsection{graph.md 中实际给出代码的算法/模板}
这里只统计 \texttt{graph.md} 中出现了具体 Python 代码的算法或可直接复用的模板；只有文字说明、没有代码实现的内容不列入“有代码算法”。例如 \texttt{Kruskal} 在 \texttt{graph.md} 中只有文字定义和图示，没有给出并查集代码，因此不放入本表的“有代码算法”清单。

\begin{longtable}{L{4.0cm}L{11.2cm}}
\toprule
\textbf{代码算法/模板} & \textbf{核心用途} \\
\midrule
Graph/Vertex 邻接表模板 & 建立邻接表图：用对象保存顶点、邻接点、边权；后续词梯、Dijkstra、Prim 都能套用。\\
Graph 邻接矩阵模板 & 建立邻接矩阵图：用二维矩阵存边，适合点数较小、需要按编号查边权的图。\\
词梯桶建图 & 建立“只差一步”的单词图：把只差一个字母的单词放入同一通配桶，再在桶内连边。\\
BFS 图遍历 & 按最少边数分层遍历：用颜色、距离、前驱和队列，从起点一层一层扩展。\\
BFS 路径回溯 & 恢复 BFS 找到的路径：沿 \texttt{pred} 从终点回溯到起点。\\
邻接矩阵 BFS & 在矩阵图上做 BFS：把 BFS 模板改成通过矩阵取邻居。\\
数轴 BFS & 在一维状态空间找最少步数：状态是一维位置，转移为 \texttt{x-1,x+1,2*x}。\\
普通迷宫 BFS & 在等代价迷宫中找最短步数：状态是 \texttt{(x,y,steps)}，二维 \texttt{visited}，四方向扩展。\\
剩余能力状态 BFS & 在带额外限制的迷宫中找最短步数：状态加入 \texttt{remaining}，判重数组变为三维。\\
不同代价迷宫 Dijkstra & 在不同格子代价的迷宫中找最小总代价：普通格代价 1，守卫格代价 2，用 \texttt{heapq} 按总代价扩展。\\
骑士周游建图 & 把棋盘问题转成图问题：把棋盘格转成图顶点，把马走日合法移动转成边。\\
骑士周游 DFS 回溯 & 搜索覆盖所有棋盘格的路径：用 DFS、颜色标记、路径栈尝试走完整个图。\\
Warnsdorff 优化 & 减少骑士周游回溯次数：优先走后续可选位置最少的点。\\
通用 DFS & 深度优先遍历整张图：用发现时间、完成时间、颜色、前驱记录搜索过程。\\
非递归 DFS & 用栈实现深度优先遍历：用显式栈模拟递归 DFS。\\
城堡 DFS 染色 & 给连通区域编号并统计面积：用 DFS 给连通房间染色，统计房间数和最大面积。\\
回溯计数 & 枚举所有可行走法：走过格子标记，递归枚举方案，返回后撤销标记。\\
算 24 回溯 & 枚举数字合并顺序和运算：每次合并两个数，递归判断能否得到 24。\\
拓扑排序 & 给有依赖关系的任务排合法顺序：统计入度，用队列反复输出入度为 0 的点。\\
AOE 关键路径 & 找工程网络中的关键活动：在 DAG 上求最早/最晚发生时间和关键活动。\\
Dijkstra 优先队列版 & 求非负权图单源最短路径：用优先队列/decreaseKey 思想维护当前最短距离。\\
Dijkstra 无优先队列版 & 求小规模非负权图单源最短路径：每轮扫描未访问点中距离最小者。\\
Floyd & 求所有点对之间的最短路径：维护 \texttt{dist} 和 \texttt{prev}，逐个中转点松弛。\\
Prim 对象图版 & 建立最小生成树：在对象图上构造 MST，使用优先队列和前驱。\\
Prim heapq 版 & 建立最小生成树：用 \texttt{heapq} 替代 decreaseKey，用 \texttt{inTree} 跳过旧状态。\\
Prim 邻接矩阵版 & 在矩阵图上建立最小生成树：用堆选择当前能连接到树的最小边。\\
\bottomrule
\end{longtable}

\subsection{作业题目与 graph.md 代码算法之间的映射}
本表只映射 PDF 作业中属于图/搜索/MST/最短路的题；树和堆题主要对应 \texttt{tree.md}，不放在本节。

\begin{longtable}{L{3.5cm}L{4.5cm}L{6.7cm}}
\toprule
\textbf{PDF 作业题目} & \textbf{对应 graph.md 有代码算法} & \textbf{本题实际采用方式} \\
\midrule
1 单词序列 & 词梯桶建图 + BFS 图遍历 & 你的代码基本沿用 \texttt{buildGraph}、\texttt{Vertex/Graph}、BFS；改为从输入列表建词表，并输出 \texttt{distance+1}。\\
2 仙岛求药 & 普通迷宫 BFS & 你的代码用 \texttt{deque}、二维 \texttt{visited}、四方向扩展；起点 \texttt{@}，终点 \texttt{*}。\\
3 变换的迷宫 & 剩余能力状态 BFS 的“状态升维”思想 & graph.md 代码第三维是剩余能力；本题实际第三维是 \texttt{time \% K}，墙在特定时间可过。\\
1 拓扑排序 & 拓扑排序 \texttt{topoSort(G)} & graph.md 用普通队列；你的代码改用 \texttt{heapq} 存入度为 0 的点，保证编号小优先，并输出 \texttt{v} 前缀。\\
2 丛林中的路 & Prim 邻接矩阵版 & 你的代码使用 \texttt{Graph\_adj\_mat} + \texttt{prim(G,0)}；字母点名转编号后写入无向邻接矩阵。\\
3 A Walk Through the Forest & Dijkstra 优先队列版 + DFS/记忆化思想 & graph.md 没有这道组合题；你的代码从顶点 2 跑 Dijkstra，再沿 \texttt{dist} 严格下降方向 DFS 计数。\\
4 打怪救公主 & 普通迷宫 BFS + 状态/资源影响后续的思想 & graph.md 没有最大血量模板；你的代码把布尔 visited 改为 \texttt{max\_blood[x][y]}，更高血量才重入队。\\
5 兔子与樱花 & Dijkstra 无优先队列版 + pred 路径恢复 & 你的代码每个查询 reset 后跑 \texttt{dijkstra\_no\_pq}，再用 \texttt{pred} 回溯并格式化输出路径；不是 Floyd。\\
56 拯救行动 & 不同代价迷宫 Dijkstra & 与 graph.md 守卫迷宫 Dijkstra 代码高度对应：\texttt{x} 代价 2，其他可走格代价 1。\\
57 连接所有点的最小费用 & Prim heapq 对象图版 & 你的代码先显式建立完全图对象图，边权为曼哈顿距离，再用 \texttt{heapq} Prim 求总费用。\\
\bottomrule
\end{longtable}

\begin{warnbox}
不要把“graph.md 里有 Floyd 代码”直接映射到“兔子与樱花”。虽然题型上 Floyd 也能做多源多查询，但你当前 PDF 和 Python 文件中的实际方法是无优先队列 Dijkstra。学案复习应以实际代码为准。
\end{warnbox}

\section{考前落地检查表}
\begin{checklistbox}
写代码前按顺序问：
\begin{enumerate}
  \item 我现在用的是 md 中哪个模板？它在本题中具体改了哪几个输入变量？
  \item PDF/同名 Python 文件里实际方法是什么？是手写结构、对象图，还是库函数？
  \item 题目答案是距离、时间、路径、数量、剩余资源，还是连通总代价？
  \item \texttt{visited} 是布尔、三维状态，还是最优值数组？
  \item 如果用了堆，堆里第一个元素代表什么：距离、边权、编号优先级，还是堆节点值？
  \item 如果用了树递归，根从哪里来：前序开头、后序结尾、当前字符、层序队首，还是编号 1？
  \item 输出口径是否和样例一致：不可达输出 0、-1、Oop!、Impossible，还是别的？
\end{enumerate}
\end{checklistbox}

\section{最该优先复习的错题组合}
\begin{enumerate}
  \item \textbf{实现堆结构 + Huffman：}一个手写 \texttt{BinHeap}，一个用 \texttt{heapq}；不要混淆。
  \item \textbf{丛林中的路 + 连接所有点：}两个都是 Prim，但前者邻接矩阵显式边，后者对象图完全图。
  \item \textbf{兔子与樱花 + A Walk Through the Forest：}两个都用了 Dijkstra 思想，但前者为每次查询找路径，后者从家出发算距离后 DFS 计数。
  \item \textbf{仙岛求药 + 变换的迷宫 + 拯救行动 + 打怪救公主：}四个都是迷宫，但分别是普通 BFS、时间模数 BFS、Dijkstra、最大资源 BFS。
  \item \textbf{22 + 53 + 50 + 23：}四个都是树输入解析，但根的位置和空节点处理完全不同。
\end{enumerate}

\end{document}
